\documentclass[10pt]{article}

\usepackage[margin=0.72in]{geometry}
\usepackage{fontspec}
\usepackage{kotex}
\usepackage{xcolor}
\usepackage{array}
\usepackage{booktabs}
\usepackage{longtable}
\usepackage{enumitem}
\usepackage{amssymb}
\usepackage{tcolorbox}
\usepackage{hyperref}

\setmainfont{Noto Serif CJK KR}
\setsansfont{Noto Sans CJK KR}
\setmonofont{DejaVu Sans Mono}

\hypersetup{
  colorlinks=true,
  linkcolor=blue!50!black,
  urlcolor=blue!50!black
}

\definecolor{TodoGray}{RGB}{105,105,105}
\definecolor{ProgressBlue}{RGB}{35,90,170}
\definecolor{DoneGreen}{RGB}{35,130,80}
\definecolor{BlockRed}{RGB}{170,55,55}
\definecolor{LightBlue}{RGB}{235,243,255}
\definecolor{LightGray}{RGB}{248,248,248}

\newcommand{\todo}{\textcolor{TodoGray}{\(\square\)}}
\newcommand{\progress}{\textcolor{ProgressBlue}{[{\raisebox{0.1ex}{\(\sim\)}}]}}
\newcommand{\done}{\textcolor{DoneGreen}{\(\boxtimes\)}}
\newcommand{\blocked}{\textcolor{BlockRed}{[!]}}
\newcommand{\code}[1]{\texttt{#1}}
\newcolumntype{L}[1]{>{\raggedright\arraybackslash}p{#1}}

\setlist[itemize]{leftmargin=1.35em, itemsep=0.16em, topsep=0.25em}
\setlist[itemize,1]{label=\todo}
\setlist[itemize,2]{label=--, leftmargin=1.8em}

\title{\vspace{-1.5em}Wind-Deformable 3DGS\\구현 현황 체크리스트}
\author{Wind3DGS}
\date{Created: 2026-05-06 \quad Last updated: 2026-05-18}

\begin{document}
\maketitle

\begin{tcolorbox}[
  colback=LightBlue,
  colframe=blue!45!black,
  title=Development Pipeline,
  fonttitle=\bfseries
]
\begin{center}
\code{3DGS input}
\(\rightarrow\)
\code{fixed-topology proxy}
\(\rightarrow\)
\code{Gaussian binding}
\(\rightarrow\)
\code{wind deformation}
\(\rightarrow\)
\code{Gaussian transport}
\(\rightarrow\)
\code{SH residual correction}
\(\rightarrow\)
\code{Gaussian rendering}
\end{center}

이 문서는 구현 진행 상황을 관리하기 위한 status board다. 연구 방향 자체가 바뀌는 결정은 기존 규칙대로 \code{ideas/idea\_sketch.tex}와 \code{ideas/changelog.md}에 기록한다.
\end{tcolorbox}

\section*{상태 표기}

\begin{center}
\begin{tabular}{@{}ll@{}}
\toprule
상태 & 의미 \\
\midrule
\todo & 아직 시작하지 않음 \\
\progress & 진행 중 \\
\done & 완료 \\
\blocked & 막힘 또는 결정 필요 \\
\bottomrule
\end{tabular}
\end{center}

\section*{실험 폴더 naming convention}

앞으로 새 구현 실험은 시간순 \code{expNNN} prefix보다 development milestone tag를 우선 사용한다.

\begin{tcolorbox}[colback=LightGray, colframe=black!20, title=권장 형식, fonttitle=\bfseries]
\begin{center}
\code{experiments/M\#\#\_short\_name/}
\end{center}

\begin{itemize}
  \item \path{experiments/M01_static_3dgs_io/}
  \item \path{experiments/M02_mesh_proxy_binding/}
  \item \path{experiments/M03_procedural_wind/}
\end{itemize}
\end{tcolorbox}

여기서 \code{M\#\#}는 이 체크리스트의 milestone 번호와 대응된다. 파일 정렬이 편하도록 folder name에는 \code{M1} 대신 \code{M01}, \code{M2} 대신 \code{M02}처럼 두 자리 숫자를 쓴다.

기존 \path{experiments/exp001_*}, \path{exp002_*}, \path{exp003_*} 폴더는 초기 planned experiment slots로 보고, 명시적으로 정리하기 전까지는 유지한다.

\section*{코드와 실험 산출물 분리 원칙}

\begin{itemize}
  \item reusable implementation은 \path{code/wind3dgs/} 아래 module로 관리한다.
  \item \path{experiments/}는 experiment README, generated assets, outputs, reports, 그리고 backward-compatible wrapper만 둔다.
  \item 기존 experiment-local 실행 명령은 wrapper로 유지한다.
  \item 새 개발은 \path{code/wind3dgs/} module 경로를 기준으로 한다.
\end{itemize}

\section*{2026-05-18 Scope Control Update}

전체 시스템은 너무 많은 구성 요소를 한 번에 핵심 기여처럼 보이지 않게 관리한다.

\begin{tcolorbox}[colback=LightBlue, colframe=blue!45!black, title=Core runtime stack, fonttitle=\bfseries]
\begin{center}
\code{proxy deformation}
\(\rightarrow\)
\code{Gaussian mean/covariance/frame transport}
\(\rightarrow\)
\code{lightweight SH residual appearance correction}
\(\rightarrow\)
\code{3DGS rendering}
\end{center}
\end{tcolorbox}

\begin{itemize}
  \item \code{SH residual}은 제거하지 않는다. 이것은 per-frame SH optimization 없이 runtime rendering을 가능하게 하는 필수 appearance correction module이다.
  \item 다만 논문 headline은 \code{AI SH prediction}이 아니라 \code{runtime wind-deformable 3DGS representation}으로 유지한다.
  \item controllable material fields, simulation anchors vs. load samples, gated floater likelihood, projection-risk-aware load sampling은 core를 보조하는 설계 장치다.
  \item MVP critical path에서는 simple fixed simulation anchors와 simple fixed material-space load samples부터 구현한다.
  \item \code{view-masked gated floater}와 \code{projection-risk-aware load sampling}은 robustness/refinement 단계로 두며, M4/M7/M8의 첫 성공 조건을 막지 않는다.
  \item advanced anchor/load scoring의 상세 기준은 \path{ideas/anchor_method_revision_notes_2026-05-11_v3.tex}를 reference로 유지한다.
\end{itemize}

\section*{현재 최우선 확인 질문}

가장 먼저 확인해야 할 기술 질문은 다음이다.

\begin{quote}
\textbf{proxy mesh가 변형될 때, 그 위에 bound된 Gaussian Splats가 position, orientation, covariance, rendering까지 자연스럽게 따라오는가?}
\end{quote}

여기에 runtime 질문을 하나 더 붙인다.

\begin{quote}
\textbf{per-frame SH optimization 없이, transported Gaussian에 lightweight SH residual correction을 붙였을 때 rendering quality와 temporal stability가 유지되는가?}
\end{quote}

첫 질문을 먼저 검증한 뒤 \code{3DGS-to-mesh extraction}으로 넘어간다. 처음부터 mesh extraction을 붙이면, 결과가 깨졌을 때 원인이 extraction noise인지, Gaussian binding 문제인지, covariance transform 문제인지 구분하기 어렵다. 두 번째 질문은 geometry transport가 안정화된 뒤 반드시 M9에서 검증한다.

\section*{Milestone Overview}

\small
\renewcommand{\arraystretch}{1.22}
\begin{longtable}{@{}L{0.065\linewidth}L{0.255\linewidth}L{0.505\linewidth}L{0.055\linewidth}@{}}
\toprule
ID & Milestone & 완료 기준 & 상태 \\
\midrule
M0 & Project and Experiment Hygiene & 첫 구현 실험 README에 goal, input, command, output, metrics, notes가 정리되어 있다. & \done \\
M1 & Static 3DGS I/O Baseline & static rendering이 안정적으로 동작하고, 이후 deformed rendering의 reference로 사용할 수 있다. & \progress \\
M2 & Synthetic Mesh-Proxy Binding MVP & mesh deformation에 따라 Gaussian motion이 안정적이며 popping, drift, covariance artifact가 없다. & \progress \\
M3 & Procedural Wind Deformation & mesh-bound Gaussian asset을 wind parameters로 scriptable하게 animation할 수 있다. & \done \\
M4 & Lightweight Mesh Simulator & proxy mesh가 단순 position warping이 아니라 connected topology와 constraints를 통해 wind에 반응한다. & \todo \\
M5 & Real 3DGS Asset to Proxy Binding & automatic mesh extraction 없이도 real 3DGS asset을 proxy mesh를 통해 animation할 수 있다. & \todo \\
M6 & 3DGS-to-Mesh Proxy Extraction & 적어도 하나의 extraction path가 stable binding과 wind simulation에 사용할 수 있는 proxy mesh를 만든다. & \todo \\
M7 & Wind Exposure From Gaussian Opacity & occluded 또는 shielded regions가 wind에 덜 반응하는 모습을 controllable하고 visually plausible하게 보여준다. & \todo \\
M8 & Force-Space Load Model & model이 stable simulation behavior를 유지하면서 wind response를 개선한다. & \todo \\
M9 & Runtime SH Residual Appearance Correction & learned residual correction이 canonical 또는 analytic-only appearance보다 quality gap을 줄이고, per-frame SH optimization 없이 runtime rendering path가 동작한다. & \todo \\
M10 & Evaluation and Paper Evidence & main paper claim을 지지할 수 있고 key ablations가 reproducible한 실험 결과가 있다. & \todo \\
\bottomrule
\end{longtable}
\normalsize

\newpage

\section{M0. Project and Experiment Hygiene}

\textbf{목표.} 첫 구현 실험을 재현 가능하고 추적 가능한 형태로 시작한다.

\begin{itemize}
  \item[\done] 첫 active implementation experiment 이름을 정한다.
    \begin{itemize}
      \item 추천: \path{experiments/M02_mesh_proxy_binding/}
    \end{itemize}
  \item[\done] 코딩 전에 experiment README를 만든다.
  \item[\done] 실행 명령어, 입력, 출력, 관찰 내용을 experiment README에 기록한다.
  \item[\done] 큰 generated outputs는 \path{code/outputs/}, \path{assets/}, 또는 experiment folder 아래에 둔다.
  \item[\done] reusable implementation code를 \path{code/wind3dgs/} package로 분리한다.
    \begin{itemize}
      \item 현재 상태: M01/M02 Python 구현은 \path{code/wind3dgs/}로 이동했고, \path{experiments/}에는 wrapper와 산출물이 남아 있다.
    \end{itemize}
\end{itemize}

\textbf{완료 기준.} 첫 구현 실험 README에 goal, input, command, output, metrics, notes가 정리되어 있다.

\section{M1. Static 3DGS I/O Baseline}

\textbf{목표.} deformation을 넣기 전에 static 3DGS asset을 정상적으로 load하고 render할 수 있는지 확인한다.

\begin{itemize}
  \item[\done] 작은 static 3DGS asset 하나를 선택하거나 준비한다.
    \begin{itemize}
      \item 현재 상태: \path{experiments/M01_static_3dgs_io/assets/synthetic_leaf_3dgs.ply}를 생성했다.
      \item asset은 leaf-like thin sheet 형태의 synthetic Inria-style 3DGS \path{.ply}다.
    \end{itemize}
  \item[\done] Gaussian attributes loader를 구현하거나 기존 코드를 연결한다.
    \begin{itemize}
      \item mean / position
      \item scale
      \item rotation
      \item opacity
      \item SH coefficients
      \item 현재 상태: ASCII / \code{binary\_little\_endian} vertex-only \path{.ply} loader를 구현했다.
      \item \code{scale\_i}는 \code{exp}, \code{opacity}는 \code{sigmoid}, \code{rot\_i}는 normalized quaternion으로 변환한다.
      \item \code{f\_dc}와 \code{f\_rest}를 SH tensor로 구성하며 기본 render는 \code{sh\_degree=0}을 사용한다.
    \end{itemize}
  \item[\done] camera intrinsics와 extrinsics를 load한다.
    \begin{itemize}
      \item 현재 상태: \path{cameras/synthetic_leaf_3dgs_cameras.json}에 \code{K}, \code{camera\_to\_world}, \code{view\_matrix}를 저장하고 load한다.
      \item convention은 \code{world-to-camera, +z forward}로 고정했다.
    \end{itemize}
  \item[\done] canonical static views를 render한다.
    \begin{itemize}
      \item 현재 상태: \code{canonical\_front}, \code{canonical\_oblique}, \code{canonical\_side} PNG를 저장했다.
    \end{itemize}
  \item[\done] 짧은 turntable 또는 fixed camera path 영상을 render한다.
    \begin{itemize}
      \item 현재 상태: 48-frame turntable PNG와 \path{outputs/turntable.gif}를 저장했다.
    \end{itemize}
  \item[\done] sanity statistics를 출력한다.
    \begin{itemize}
      \item Gaussian count
      \item bounding box
      \item opacity range
      \item scale range
      \item SH coefficient shape
      \item 현재 상태: \path{outputs/stats.json}와 \path{outputs/render_report.md}에 기록한다.
    \end{itemize}
  \item[\done] baseline images 또는 video를 experiment/output folder에 저장한다.
    \begin{itemize}
      \item 현재 상태: \path{outputs/} 아래에 canonical renders와 turntable preview를 저장했다.
      \item 사용자 local shell에서 \code{--backend cpu\_debug} 기준 canonical renders 3장과 turntable frames 48장 생성이 확인되었다.
    \end{itemize}
  \item[\blocked] \code{gsplat} CUDA Gaussian rasterizer로 static rendering을 검증한다.
    \begin{itemize}
      \item 현재 상태: \path{render_static_baseline.py}에 \code{auto}, \code{gsplat}, \code{cpu\_debug} backend를 구현했다.
      \item 1차 blocker였던 CUDA Toolkit / \code{nvcc} 누락은 사용자 local shell에서 해결했다.
      \item 확인 결과: local GPU는 \code{NVIDIA GeForce GTX 1080 Ti}, Compute Capability \code{6.1}이다.
      \item 현재 blocker: \code{gsplat 1.5.3} JIT build가 \code{compute\_61} target에서 CUDA cooperative groups compile error로 실패한다.
      \item 원인: NVIDIA CUDA docs 기준 \code{labeled\_partition}은 Compute Capability 7.0 이상이 필요하다. GTX 1080 Ti는 CC 6.1이라 현재 \code{gsplat 1.5.3} CUDA backend와 맞지 않는다.
      \item 다음 액션: 이 machine에서는 M1을 \code{cpu\_debug} fallback으로 유지하고, 논문용 \code{gsplat} render 검증은 CC >= 7.0 GPU에서 진행한다.
    \end{itemize}
  \item ground-truth views가 있으면 PSNR, SSIM, LPIPS를 계산한다.
    \begin{itemize}
      \item 현재 상태: synthetic/minimal asset에는 ground-truth image가 없어 보류한다.
    \end{itemize}
\end{itemize}

\textbf{완료 기준.} synthetic/minimal static pipeline은 현재 machine에서 \code{cpu\_debug} 기준으로 재현 가능하다. 다만 GTX 1080 Ti에서는 \code{gsplat 1.5.3} CUDA backend가 hardware capability 때문에 막혀 있으므로, 논문 품질의 static rendering reference는 CC >= 7.0 GPU에서 추가 확인해야 한다.

\section{M2. Synthetic Mesh-Proxy Binding MVP}

\textbf{목표.} 깨끗한 synthetic mesh deformation을 기준으로 Gaussian Splats가 제대로 따라 움직이는지 검증한다.

\begin{tcolorbox}[colback=LightGray, colframe=black!20, title=mesh extraction을 뒤로 미루는 이유, fonttitle=\bfseries]
이 단계에서는 일부러 \code{3DGS-to-mesh extraction}을 하지 않는다. 먼저 깨끗한 synthetic mesh를 사용해서 Gaussian binding과 transform bug를 분리해서 본다.
\end{tcolorbox}

\begin{itemize}
  \item[\done] 간단한 proxy mesh를 생성하거나 load한다.
    \begin{itemize}
      \item flag plane
      \item cloth strip
      \item leaf-like thin sheet
    \end{itemize}
  \item[\done] mesh surface 위 또는 근처에 synthetic Gaussians를 sample한다.
  \item[\done] M1 synthetic Inria-style PLY를 M2 viewer에서 load하고 generated proxy mesh에 bind한다.
    \begin{itemize}
      \item 현재 상태: \path{viewer_gpu.py}에 \code{--ply} 로드 경로를 추가했다.
      \item 예시 asset: \path{experiments/M01_static_3dgs_io/assets/synthetic_leaf_3dgs.ply}
      \item 현재 기본 방식: PLY Gaussian mean을 proxy grid에 rasterize하고, Gaussian이 존재하는 cell만 남기는 occupancy-extracted proxy mesh를 생성한다.
      \item \code{--ply-proxy-mode bbox}를 사용하면 이전 rectangular bounding-box proxy로 비교할 수 있다.
      \item 각 Gaussian은 proxy triangle의 \code{triangle\_id}, \code{barycentric coordinates}, \code{local offset}으로 bind된다.
      \item PLY의 \code{rot\_0..3} quaternion은 anisotropic frame으로 변환한 뒤 triangle-local basis를 통해 transport한다.
    \end{itemize}
  \item[\done] 각 Gaussian을 mesh triangle에 bind한다.
    \begin{itemize}
      \item triangle id
      \item barycentric coordinates
      \item triangle local frame
      \item local offset
    \end{itemize}
  \item[\done] 단순하고 원인을 알 수 있는 mesh deformations를 적용한다.
    \begin{itemize}
      \item bending
      \item sine-wave flutter
      \item twist
      \item edge flap
      \item compound
      \item 현재 상태: \path{viewer_gpu.py}와 \path{viewer.html} 모두 UI에서 \code{sine}, \code{bend}, \code{twist}, \code{edge\_flap}, \code{compound}를 선택할 수 있다.
    \end{itemize}
  \item[\done] deformed triangle의 barycentric position으로 Gaussian centers를 update한다.
  \item[\done] triangle-local rotation으로 Gaussian covariance 또는 anisotropic frame을 transport한다.
    \begin{itemize}
      \item 현재 상태: canonical anisotropic frame coefficients를 triangle-local coordinates로 보존하고, deformed triangle basis에서 frame과 covariance를 재구성한다.
      \item \path{viewer_gpu.py}: sampled transported GS ellipsoid surface를 기본으로 표시하고, 필요하면 frame axes도 표시한다.
      \item \path{viewer.html}: anisotropic ellipse와 frame axes를 표시한다.
    \end{itemize}
  \item[\progress] frame flip, seam popping, exploding covariance, surface sliding이 생기는지 확인한다.
    \begin{itemize}
      \item 현재 상태: visual inspection과 numeric binding/covariance verification은 통과. \code{10x10}, \code{30x30}, \code{50x50} asset에 대해 \code{sine}, \code{bend}, \code{twist}, \code{edge\_flap}, \code{compound} deformation을 모두 검증했다.
      \item 실제 3DGS rasterization과 SH appearance artifact 검증은 아직 전이다.
    \end{itemize}
  \item[\done] dense visual inspection을 위한 GPU-backed debug viewer를 추가한다.
    \begin{itemize}
      \item 현재 상태: \path{viewer_gpu.py}는 \code{ModernGL + GLFW} 기반이며, \code{10x10}, \code{30x30}, \code{50x50} smoke test를 통과했다.
      \item \code{tkinter} control panel로 web viewer의 asset selector, mesh statistics, display toggles, deformation sliders, reset view, loaded file path UI를 제공한다.
      \item per-frame vertex normal shading과 distance-based depth cueing을 추가해 cloth 굴곡과 앞뒤 depth 구분을 개선했다.
      \item face surface는 opaque로 표시한다.
      \item back face culling은 기본 on이며, camera side에 따라 front-face winding을 \code{ccw}/\code{cw}로 보정한다. UI toggle / \code{C} key / \code{--no-backface-culling}으로 끌 수 있다.
      \item camera rotation은 yaw/pitch/roll 3축을 지원한다. left drag는 yaw/pitch, right drag 또는 \code{Shift}+left drag는 roll, UI slider와 arrow/\code{Z}/\code{X} key도 지원한다.
      \item transported GS ellipsoid surface는 UI toggle / \code{E} key / \code{--no-gaussian-ellipsoids}로 켜고 끌 수 있다.
      \item transported GS anisotropic frame axes는 UI toggle / \code{O} key / \code{--show-gaussian-frames}로 켜고 끌 수 있다.
      \item \code{--ply} CLI와 Tk control panel의 \code{load PLY...} / \code{use sample GS} 버튼으로 M1 synthetic PLY와 built-in sample GS를 전환할 수 있다.
      \item PLY mode는 \code{--ply-proxy-mode occupancy}를 기본으로 사용해 leaf 영역 밖 rectangular mesh를 줄인다.
    \end{itemize}
  \item[\progress] before/after frames 또는 짧은 animation을 render한다.
    \begin{itemize}
      \item 현재 상태: \path{viewer.html}은 zero-setup fallback이고, \path{viewer_gpu.py}는 dense sample Gaussian inspection 용도다. 실제 3DGS rendering은 아직 전이다.
    \end{itemize}
\end{itemize}

\textbf{완료 기준.} mesh deformation에 따라 Gaussian motion이 안정적이며, 눈에 띄는 popping, drift, covariance artifact가 없다.

\section{M3. Procedural Wind Deformation}

\textbf{목표.} manual deformation을 단순한 wind-driven procedural deformation으로 바꾼다.

\begin{itemize}
  \item[\done] lightweight wind field를 정의한다.
    \begin{itemize}
      \item global direction
      \item strength
      \item gust frequency
      \item spatial phase
      \item 현재 상태: \path{code/wind3dgs/m03_procedural_wind/wind_field.py}에 one-way \code{WindParameters}와 procedural gust deformation을 추가했다.
      \item M3 wind는 global wind direction + procedural spatial gust 방식이며, voxel wind grid나 two-way coupling은 아직 사용하지 않는다.
    \end{itemize}
  \item[\done] synthetic mesh에 wind-like deformation을 적용한다.
    \begin{itemize}
      \item 현재 상태: \path{viewer_gpu.py}의 \code{--deformation wind}로 M02 cloth proxy와 M1 synthetic leaf PLY occupancy proxy에 같은 wind deformation path를 적용할 수 있다.
      \item anchored boundary는 고정되고, free vertices만 wind strength, gust frequency, spatial scale, turbulence에 반응한다.
    \end{itemize}
  \item[\done] artist-style controls를 지원한다.
    \begin{itemize}
      \item wind direction
      \item amplitude
      \item frequency
      \item anchored boundary
      \item 현재 상태: Tk control panel과 CLI에 \code{wind direction}, \code{wind spatial scale}, \code{turbulence}를 추가했다.
      \item 기존 \code{amplitude}는 wind strength, \code{frequency}는 gust frequency로 사용한다.
    \end{itemize}
  \item[\done] 다음 조건을 비교한다.
    \begin{itemize}
      \item position-only Gaussian transport
      \item position plus covariance transport
      \item optional local-frame appearance correction
      \item 현재 상태: \code{--transport-mode full}은 Gaussian center와 anisotropic frame/covariance를 함께 transport한다.
      \item 현재 상태: \code{--transport-mode position\_only}는 center만 mesh에 bind하고 anisotropic frame은 canonical frame으로 유지해 비교 baseline으로 쓴다.
      \item local-frame / SH residual appearance correction은 M9의 core runtime module로 남긴다. M3에서는 geometry transport path만 검증한다.
    \end{itemize}
  \item[\done] 여러 wind settings에 대한 qualitative videos를 저장한다.
    \begin{itemize}
      \item 현재 상태: \path{code/wind3dgs/m03_procedural_wind/render_wind_preview.py}를 추가했다.
      \item OpenGL viewer 없이 proxy mesh와 bound Gaussian centers를 orthographic CPU preview로 렌더링해 GIF, poster PNG, JSON report를 저장한다.
      \item \code{calm}, \code{crosswind}, \code{gusty} preset을 \code{50x50} cloth sample GS와 M1 synthetic leaf PLY occupancy proxy에 대해 생성했다.
    \end{itemize}
\end{itemize}

\textbf{완료 기준.} mesh-bound Gaussian asset을 wind parameters로 scriptable하게 animation할 수 있다. 현재 module command와 experiment wrapper smoke test 기준으로 scriptable wind animation path가 동작하고, headless qualitative GIF도 생성된다.

\section{M4. Lightweight Mesh Simulator}

\textbf{목표.} procedural deformation에서 topology-aware simulator로 넘어간다.

\begin{itemize}
  \item 작은 mass-spring, cloth, 또는 XPBD-style solver를 구현하거나 연결한다.
  \item persistent simulation anchors를 정의한다.
    \begin{itemize}
      \item 기본 MVP: fixed proxy mesh vertices를 simulation anchors로 사용한다.
      \item optional: 나중에 reduced simulation control nodes를 별도로 sample한다.
    \end{itemize}
  \item anchored vertices 또는 boundary constraints를 추가한다.
  \item stiffness, damping, mass parameters를 추가한다.
  \item relative velocity와 surface normal을 사용한 wind force를 추가한다.
  \item simple fixed material-space load samples를 정의한다.
    \begin{itemize}
      \item MVP에서는 uniform face samples 또는 face centers부터 시작한다.
      \item \code{projection-risk-aware} refinement는 M8 이후 optional로 둔다.
    \end{itemize}
  \item flag/cloth-strip test cases에서 안정적으로 simulation을 돌린다.
  \item Gaussian transport를 위해 mesh states over time을 export한다.
\end{itemize}

\textbf{완료 기준.} proxy mesh가 단순 position warping이 아니라 connected topology와 constraints를 통해 wind에 반응한다.

\section{M5. Real 3DGS Asset to Proxy Binding}

\textbf{목표.} binding pipeline을 real 또는 pretrained 3DGS assets에 연결한다.

\begin{itemize}
  \item M1 loader를 통해 real/pretrained 3DGS asset을 load한다.
  \item low-opacity floaters를 제거하거나 mask한다.
  \item 필요하면 surface-oriented Gaussian attributes를 추정한다.
  \item 처음에는 manually provided proxy mesh 또는 simple proxy mesh로 시작한다.
  \item real Gaussians를 proxy mesh에 bind한다.
  \item real asset에 대해 M2/M3 deformation path를 실행한다.
  \item static result와 animated result를 render한다.
\end{itemize}

\textbf{완료 기준.} automatic mesh extraction 없이도 real 3DGS asset을 proxy mesh를 통해 animation할 수 있다.

\section{M6. 3DGS-to-Mesh Proxy Extraction}

\textbf{목표.} manually provided proxy mesh를 extracted simulation proxy로 대체한다.

\begin{itemize}
  \item candidate extraction paths를 테스트한다.
    \begin{itemize}
      \item SuGaR-style mesh extraction
      \item Gaussian Opacity Fields-style extraction
      \item TSDF fusion from rendered depths
      \item point-cloud Poisson reconstruction
    \end{itemize}
  \item extracted geometry를 simulation용으로 simplify 또는 remesh한다.
  \item proxy quality를 평가한다.
    \begin{itemize}
      \item topology stability
      \item triangle quality
      \item thin-structure preservation
      \item reliability-weighted Gaussian transport demand
      \item fixed-topology simulation anchorability
    \end{itemize}
  \item optional robustness score를 추가한다.
    \begin{itemize}
      \item raw Gaussian count만 proxy 중요도로 쓰지 않는다.
      \item opacity, surface distance, optional view evidence를 사용해 \code{R\_G}를 계산한다.
      \item direct isolation reliability 대신 \code{view-masked gated floater likelihood}를 사용한다.
      \item 이 항목은 M6 첫 extraction success의 blocker가 아니라 v3-based refinement다.
    \end{itemize}
  \item render-quality mesh extraction과 simulation-oriented proxy extraction을 비교한다.
\end{itemize}

\textbf{완료 기준.} 적어도 하나의 extraction path가 stable binding과 wind simulation에 사용할 수 있는 proxy mesh를 만든다.

\section{M7. Wind Exposure From Gaussian Opacity}

\textbf{목표.} object의 어느 부분이 wind에 노출되거나 가려지는지 추정한다.

\begin{itemize}
  \item 현재 Gaussian state를 wind axis 방향으로 rasterize한다.
  \item approximate transmittance 또는 exposure map을 계산한다.
  \item fixed material-space load sample에서 exposure를 sample한다.
    \begin{itemize}
      \item 논문 notation: \(q_\ell := (\tau_\ell,\mathbf{b}_\ell)\)
      \item \(\tau_\ell\): proxy face/triangle index
      \item \(\mathbf{b}_\ell\): 해당 face 위 barycentric weights
      \item frame마다 anchor를 다시 고르지 않는다.
      \item load sample identity는 preprocessing 이후 고정한다.
    \end{itemize}
  \item load-sample force를 barycentric accumulation으로 persistent simulation anchors에 project한다.
  \item exposure를 사용해 wind force를 modulate한다.
  \item with exposure와 without exposure animation을 비교한다.
\end{itemize}

\textbf{완료 기준.} occluded 또는 shielded regions가 wind에 덜 반응하는 모습을 controllable하고 visually plausible하게 보여준다.

\section{M8. Force-Space Load Model}

\textbf{목표.} physical integration을 명시적으로 유지하면서 aerodynamic loads를 learn 또는 calibrate한다.

\begin{itemize}
  \item 각 load sample과 simulation anchor의 local state features를 정의한다.
  \item direct deformation이 아니라 forces 또는 torques를 predict한다.
  \item analytic 또는 parameterized load model을 먼저 구현한다.
  \item dense simulation labels가 있으면 dense forces를 sparse proxy anchors로 project한다.
    \begin{itemize}
      \item 이 항목은 core MVP blocker가 아니라 supervised refinement다.
    \end{itemize}
  \item net force와 wrench에 대한 conservation checks를 추가한다.
  \item optional: dense-force correction ratio로 \code{projection-risk-aware} load-sample refinement를 계산한다.
  \item minimal MLP 또는 learned residual load model을 train한다.
  \item analytic-only loads와 learned residual loads를 비교한다.
\end{itemize}

\textbf{완료 기준.} model이 stable simulation behavior를 유지하면서 wind response를 개선한다.

\section{M9. Runtime SH Residual Appearance Correction}

\textbf{목표.} expensive per-frame SH optimization 없이 deformation 이후 view-dependent appearance를 보존한다.

이 milestone은 optional polish가 아니라 runtime system의 필수 appearance stage다. 단, 논문 headline은 \code{AI SH prediction}이 아니라 proxy deformation, Gaussian transport, SH residual runtime correction의 결합으로 유지한다.

\begin{itemize}
  \item appearance baselines를 세운다.
    \begin{itemize}
      \item canonical SH only
      \item analytic local-frame correction
      \item per-frame SH optimization upper bound
      \item learned SH residual
    \end{itemize}
  \item residual labels 또는 pseudo-labels를 정의한다.
  \item local deformation과 wind features를 사용해 lightweight residual predictor를 train한다.
    \begin{itemize}
      \item 입력 후보: local frame rotation, normal change, stretch, curvature change, wind direction/strength, material latent, canonical SH, binding features
    \end{itemize}
  \item runtime inference path에 residual predictor를 연결한다.
  \item per-frame SH optimization 없이 animation sequence를 render한다.
  \item rendering quality와 temporal flicker를 평가한다.
\end{itemize}

\textbf{완료 기준.} learned residual correction이 canonical 또는 analytic-only appearance보다 의미 있는 quality gap을 줄이고, per-frame SH optimization 없이 runtime rendering path가 동작한다.

\section{M10. Evaluation and Paper Evidence}

\textbf{목표.} 구현 결과를 paper claim을 뒷받침하는 evidence로 만든다.

\begin{itemize}
  \item synthetic scenes를 준비한다.
    \begin{itemize}
      \item flag
      \item cloth strip
      \item curtain
      \item leaf or plant-like sheet
    \end{itemize}
  \item 가능하면 real 3DGS scenes를 준비한다.
  \item main baselines를 실행한다.
    \begin{itemize}
      \item static 3DGS
      \item direct Gaussian deformation
      \item mesh-proxy deformation
      \item with and without exposure
      \item with and without SH residual
      \item canonical SH vs analytic SH correction vs learned SH residual
    \end{itemize}
  \item metrics를 측정한다.
    \begin{itemize}
      \item PSNR / SSIM / LPIPS
      \item temporal consistency
      \item trajectory or surface displacement error
      \item runtime / FPS
      \item memory
    \end{itemize}
  \item qualitative videos와 failure cases를 저장한다.
  \item 연구 방향이 바뀌면 \code{ideas/idea\_sketch.tex}와 \code{ideas/changelog.md}를 업데이트한다.
\end{itemize}

\textbf{완료 기준.} main paper claim을 지지할 수 있고 key ablations가 reproducible한 실험 결과가 있다.

\section*{Open Decisions}

\begin{itemize}
  \item[\done] base로 사용할 3DGS implementation 또는 renderer는 무엇인가?
    \begin{itemize}
      \item 결정: M1/M5의 우선 renderer는 \code{gsplat} CUDA Gaussian rasterizer로 두고, 개발 중에는 \code{cpu\_debug} fallback을 유지한다.
    \end{itemize}
  \item[\done] M1의 첫 static asset은 무엇으로 할 것인가?
    \begin{itemize}
      \item 결정: real/pretrained asset 확보 전까지 \path{M01_static_3dgs_io}의 synthetic leaf-like Inria-style 3DGS asset을 사용한다.
    \end{itemize}
  \item[\done] M2는 custom minimal renderer, existing 3DGS renderer, 또는 둘 다 사용할 것인가?
    \begin{itemize}
      \item 결정: M02는 custom minimal GPU debug renderer를 사용하고, \path{viewer.html}은 fallback으로 유지한다. existing 3DGS renderer 연결은 M1/M5 이후로 미룬다.
    \end{itemize}
  \item[\done] SH residual을 scope에서 뺄 것인가?
    \begin{itemize}
      \item 결정: 빼지 않는다. SH residual은 runtime에서 per-frame SH optimization을 피하기 위한 필수 appearance correction module이다.
      \item 단, headline contribution은 \code{AI SH prediction}이 아니라 runtime wind-deformable 3DGS representation으로 유지한다.
    \end{itemize}
  \item[\done] advanced anchor/load scoring을 MVP critical path에 넣을 것인가?
    \begin{itemize}
      \item 결정: 넣지 않는다. MVP는 fixed simulation anchors와 simple fixed material-space load samples부터 시작한다.
      \item \code{view-masked gated floater}와 \code{projection-risk-aware load sampling}은 robustness/refinement 단계로 둔다.
    \end{itemize}
  \item M4에서 사용할 mesh simulation library는 무엇인가?
  \item M6에서 가장 먼저 시도할 mesh extraction baseline은 무엇인가?
  \item 기존 \path{experiments/exp001_*}--\path{exp003_*} planned folders를 그대로 둘지, 나중에 milestone tag 기반으로 정리할지 결정한다.
\end{itemize}

\section*{Notes}

\begin{itemize}
  \item SH residual prediction은 M2/M3의 첫 geometry 목표는 아니지만, 최종 runtime system에서는 필수 module이다. M9에서 반드시 구현하고 M10 evidence에 포함한다.
  \item automatic mesh extraction부터 시작하지 않는다. 먼저 Gaussian-to-mesh binding과 deformation transport가 sound한지 증명한다.
  \item advanced anchor/load scoring부터 시작하지 않는다. fixed simulation anchors와 simple fixed load samples로 먼저 solver와 rendering loop를 닫는다.
  \item 각 milestone은 다음 layer를 붙이기 전에 runnable하고 visually inspectable해야 한다.
  \item M02의 GPU viewer는 구현 편의상 \code{ModernGL + GLFW}를 사용하지만, 연구 기여가 OpenGL 자체에 의존한다는 뜻은 아니다.
\end{itemize}

\end{document}
